% GENERATED FILE. Edit canonical lessons, then run npm run generate:books.
% canonical-source-hash: fnv1a64:8b2072df3c503eb6
% canonical-lessons: MW-C149-puro, MW-C149-khus, MW-C149-darpok, MW-C149-bahadur, MW-C149-gailo, MW-R149-first-pass-animals-and-describing-words, MW-R149-second-pass-describing-people

\chapter{Complete, Happy, Timid, Brave, Crazy}
\label{ch:mw-complete-happy-timid-brave-crazy}
\hlchaptermodality{149}

\begin{chapteropening}
\textbf{By the end of this chapter:} \emph{I can say complete, happy, timid, brave and crazy.}

Complete the last lesson of chapter 149: I can say complete, happy, timid, brave and crazy.
\end{chapteropening}

\section[\texorpdfstring{\mw{पूरो} (pūro)}{पूरो (pūro)}]{\mw{पूरो} (pūro) — complete, whole}

\label{lesson:MW-C149-puro}



\begin{quote}
Before the new one: say the Marwadi for plain, then the Marwadi for unusual.
\end{quote}

\subsection*{You'll want to know: \mw{पूरो}}
\textbf{\mw{पूरो}} — \emph{pūro} — \textquotedblleft{}complete, whole\textquotedblright{}.

\begin{cousinweb}
\textbf{\mw{पूरो}} (\emph{pūro}) is a tadbhava, a word inherited from Sanskrit and reshaped by sound change. It comes from \emph{pūra}, \textquotedblleft{}filling\textquotedblright{}. CDIAL lists Old Marwari \emph{puraü} under it (CDIAL 8330).
\end{cousinweb}

\begin{grammarlens}[title={What you've built}]
One more word for this chapter.
\end{grammarlens}

\subsection*{Your turn}
\begin{itemize}
\raggedright
  \item \emph{Say it:} \emph{pūro}
  \item \emph{Say it:} \emph{pūro}, once more
  \item \emph{Say it:} say \emph{pūro}
  \item \emph{Recall:} say the Marwadi for plain, then the Marwadi for unusual, then say \emph{pūro} again
\end{itemize}

\begin{tcolorbox}[breakable,colback=teal!4,colframe=teal!35!black,title={Before you move on}]
What does \mw{पूरो} mean? (\textquotedblleft{}Complete, whole\textquotedblright{}.) Say it once more.
\end{tcolorbox}
\section[\texorpdfstring{\mw{खुस} (khus)}{खुस (khus)}]{\mw{खुस} (khus) — happy}

\label{lesson:MW-C149-khus}



\begin{quote}
Before the new one: say the Marwadi for unusual, then the Marwadi for complete.
\end{quote}

\subsection*{You'll want to know: \mw{खुस}}
\textbf{\mw{खुस}} — \emph{khus} — \textquotedblleft{}happy\textquotedblright{}.

\begin{cousinweb}
\textbf{\mw{खुस}} (\emph{khus}) is a loan, usually traced to Persian \emph{xuš}, \textquotedblleft{}happy, pleasant\textquotedblright{}. Marwadi says the \emph{š} as \emph{s}.
\end{cousinweb}

\begin{grammarlens}[title={What you've built}]
One more word for this chapter.
\end{grammarlens}

\subsection*{Your turn}
\begin{itemize}
\raggedright
  \item \emph{Say it:} \emph{khus}
  \item \emph{Say it:} \emph{khus}, once more
  \item \emph{Say it:} say \emph{khus}
  \item \emph{Recall:} say the Marwadi for unusual, then the Marwadi for complete, then say \emph{khus} again
\end{itemize}

\begin{tcolorbox}[breakable,colback=teal!4,colframe=teal!35!black,title={Before you move on}]
What does \mw{खुस} mean? (\textquotedblleft{}Happy\textquotedblright{}.) Say it once more.
\end{tcolorbox}
\section[\texorpdfstring{\mw{डरपोक} (ḍarpok)}{डरपोक (ḍarpok)}]{\mw{डरपोक} (ḍarpok) — timid}

\label{lesson:MW-C149-darpok}



\begin{quote}
Before the new one: say the Marwadi for complete, then the Marwadi for happy.
\end{quote}

\subsection*{You'll want to know: \mw{डरपोक}}
\textbf{\mw{डरपोक}} — \emph{ḍarpok} — \textquotedblleft{}timid\textquotedblright{}.

\begin{cousinweb}
\textbf{\mw{डरपोक}} (\emph{ḍarpok}) is built on \emph{ḍar}, \textquotedblleft{}fear\textquotedblright{}, which CDIAL lists as Marwari \emph{ḍar} under Sanskrit \emph{dara}, \textquotedblleft{}fear\textquotedblright{} (CDIAL 6186). CDIAL has no entry for the ending \emph{-pok}, so that part is not settled.
\end{cousinweb}

\begin{grammarlens}[title={What you've built}]
One more word for this chapter.
\end{grammarlens}

\subsection*{Your turn}
\begin{itemize}
\raggedright
  \item \emph{Say it:} \emph{ḍarpok}
  \item \emph{Say it:} \emph{ḍarpok}, once more
  \item \emph{Say it:} say \emph{ḍarpok}
  \item \emph{Recall:} say the Marwadi for complete, then the Marwadi for happy, then say \emph{ḍarpok} again
\end{itemize}

\begin{tcolorbox}[breakable,colback=teal!4,colframe=teal!35!black,title={Before you move on}]
What does \mw{डरपोक} mean? (\textquotedblleft{}Timid\textquotedblright{}.) Say it once more.
\end{tcolorbox}
\section[\texorpdfstring{\mw{बहादुर} (bahādur)}{बहादुर (bahādur)}]{\mw{बहादुर} (bahādur) — brave}

\label{lesson:MW-C149-bahadur}



\begin{quote}
Before the new one: say the Marwadi for happy, then the Marwadi for timid.
\end{quote}

\subsection*{You'll want to know: \mw{बहादुर}}
\textbf{\mw{बहादुर}} — \emph{bahādur} — \textquotedblleft{}brave\textquotedblright{}.

\begin{cousinweb}
\textbf{\mw{बहादुर}} (\emph{bahādur}) is a loan, usually traced to Persian \emph{bahādur}, \textquotedblleft{}brave, hero\textquotedblright{}. Persian took it from a Turkic or Mongol word for \textquotedblleft{}hero\textquotedblright{}.
\end{cousinweb}

\begin{grammarlens}[title={What you've built}]
One more word for this chapter.
\end{grammarlens}

\subsection*{Your turn}
\begin{itemize}
\raggedright
  \item \emph{Say it:} \emph{bahādur}
  \item \emph{Say it:} \emph{bahādur}, once more
  \item \emph{Say it:} say \emph{bahādur}
  \item \emph{Recall:} say the Marwadi for happy, then the Marwadi for timid, then say \emph{bahādur} again
\end{itemize}

\begin{tcolorbox}[breakable,colback=teal!4,colframe=teal!35!black,title={Before you move on}]
What does \mw{बहादुर} mean? (\textquotedblleft{}Brave\textquotedblright{}.) Say it once more.
\end{tcolorbox}
\section[\texorpdfstring{\mw{गैलो} (gailo)}{गैलो (gailo)}]{\mw{गैलो} (gailo) — crazy}

\label{lesson:MW-C149-gailo}



\begin{quote}
Before the new one: say the Marwadi for timid, then the Marwadi for brave.
\end{quote}

\subsection*{You'll want to know: \mw{गैलो}}
\textbf{\mw{गैलो}} — \emph{gailo} — \textquotedblleft{}crazy\textquotedblright{}.

\begin{cousinweb}
\textbf{\mw{गैलो}} (\emph{gailo}) is a tadbhava, a word inherited from Sanskrit and reshaped by sound change. It comes from \emph{grahila}, \textquotedblleft{}possessed by an evil spirit\textquotedblright{}. CDIAL lists Marwari \emph{gahalo} under it (CDIAL 4366). The Marwari row means \textquotedblleft{}foolish\textquotedblright{}; \textbf{\mw{गैलो}} (\emph{gailo}) has lost its \emph{h}. CDIAL thinks the Sanskrit word was made from Middle Indo-Aryan \emph{gahila}, built on \emph{graha}, \textquotedblleft{}seizing\textquotedblright{}.
\end{cousinweb}

\begin{grammarlens}[title={What you've built}]
One more word for this chapter.
\end{grammarlens}

\subsection*{Your turn}
\begin{itemize}
\raggedright
  \item \emph{Say it:} \emph{gailo}
  \item \emph{Say it:} \emph{gailo}, once more
  \item \emph{Say it:} say \emph{gailo}
  \item \emph{Recall:} say the Marwadi for timid, then the Marwadi for brave, then say \emph{gailo} again
\end{itemize}

\begin{tcolorbox}[breakable,colback=teal!4,colframe=teal!35!black,title={Before you move on}]
What does \mw{गैलो} mean? (\textquotedblleft{}Crazy\textquotedblright{}.) Say it once more.
\end{tcolorbox}
\section[(dialogue)]{Animals and describing words — a first pass}

\label{lesson:MW-R149-first-pass-animals-and-describing-words}



\begin{quote}
Say the words for brave and crazy.
\end{quote}

\subsection*{Your turn}
Say these aloud:

\begin{itemize}
\raggedright
  \item a mouse, a fly, a mosquito, an ant, a scorpion
  \item good, bad, narrow, wide, deep
  \item light, hard, soft, wet, dry
  \item hot, salty, sour, bitter, spicy
  \item beautiful, good-looking, young, thin, short
\end{itemize}

\begin{tcolorbox}[breakable,colback=teal!4,colframe=teal!35!black,title={Before you move on}]
A mouse? (\textbf{\mw{उंदरो}}, \emph{undro}.) A fly? (\textbf{\mw{माखी}}, \emph{mākhī}.) A mosquito? (\textbf{\mw{मच्छर}}, \emph{machchhar}.) An ant? (\textbf{\mw{कीड़ी}}, \emph{kīṛī}.) A scorpion? (\textbf{\mw{बीछू}}, \emph{bīchhū}.) Good? (\textbf{\mw{चोखो}}, \emph{chokho}.) Bad? (\textbf{\mw{खराब}}, \emph{kharāb}.) Narrow? (\textbf{\mw{सांकड़ो}}, \emph{sā̃kṛo}.) Wide? (\textbf{\mw{चौड़ो}}, \emph{chauṛo}.) Deep? (\textbf{\mw{गैरो}}, \emph{gairo}.) Light? (\textbf{\mw{हळको}}, \emph{haḷko}.) Hard? (\textbf{\mw{कड़क}}, \emph{kaṛak}.) Soft? (\textbf{\mw{नरम}}, \emph{naram}.) Wet? (\textbf{\mw{गीलो}}, \emph{gīlo}.) Dry? (\textbf{\mw{सूको}}, \emph{sūko}.) Hot? (\textbf{\mw{तातो}}, \emph{tāto}.) Salty? (\textbf{\mw{खारो}}, \emph{khāro}.) Sour? (\textbf{\mw{खाटो}}, \emph{khāṭo}.) Bitter? (\textbf{\mw{कड़वो}}, \emph{kaṛvo}.) Spicy? (\textbf{\mw{तीखो}}, \emph{tīkho}.) Beautiful? (\textbf{\mw{सुंदर}}, \emph{sundar}.) Good-looking? (\textbf{\mw{रूपाळो}}, \emph{rūpāḷo}.) Young? (\textbf{\mw{जवान}}, \emph{javān}.) Thin? (\textbf{\mw{दुबळो}}, \emph{dubaḷo}.) Short? (\textbf{\mw{ठिगणो}}, \emph{ṭhigṇo}.)
\end{tcolorbox}
\section[(dialogue)]{Describing people — a second pass}

\label{lesson:MW-R149-second-pass-describing-people}



\begin{quote}
Say the words for pink and saffron-coloured.
\end{quote}

\subsection*{Your turn}
Say these aloud:

\begin{itemize}
\raggedright
  \item pink, saffron-coloured, sky-blue, purple, orange
  \item shiny, calm, clever, silly, lazy
  \item hard-working, honest, kind, miserly, good-natured
  \item wicked, necessary, special, plain, unusual
  \item complete, happy, timid, brave, crazy
\end{itemize}

\begin{tcolorbox}[breakable,colback=teal!4,colframe=teal!35!black,title={Before you move on}]
Pink? (\textbf{\mw{गुलाबी}}, \emph{gulābī}.) Saffron-coloured? (\textbf{\mw{केसरियो}}, \emph{kesariyo}.) Sky-blue? (\textbf{\mw{आसमानी}}, \emph{āsmānī}.) Purple? (\textbf{\mw{बैंगणी}}, \emph{baiṅgṇī}.) Orange? (\textbf{\mw{नारंगी}}, \emph{nāraṅgī}.) Shiny? (\textbf{\mw{चमकीलो}}, \emph{chamkīlo}.) Calm? (\textbf{\mw{सांत}}, \emph{sānt}.) Clever? (\textbf{\mw{हुसियार}}, \emph{husiyār}.) Silly? (\textbf{\mw{बावळो}}, \emph{bāvaḷo}.) Lazy? (\textbf{\mw{आळसी}}, \emph{āḷsī}.) Hard-working? (\textbf{\mw{मेहनती}}, \emph{mehantī}.) Honest? (\textbf{\mw{ईमानदार}}, \emph{īmāndār}.) Kind? (\textbf{\mw{दयाळु}}, \emph{dayāḷu}.) Miserly? (\textbf{\mw{कंजूस}}, \emph{kañjūs}.) Good-natured? (\textbf{\mw{भलो}}, \emph{bhalo}.) Wicked? (\textbf{\mw{बुरो}}, \emph{buro}.) Necessary? (\textbf{\mw{जरूरी}}, \emph{jarūrī}.) Special? (\textbf{\mw{खास}}, \emph{khās}.) Plain? (\textbf{\mw{सादो}}, \emph{sādo}.) Unusual? (\textbf{\mw{अनोखो}}, \emph{anokho}.) Complete? (\textbf{\mw{पूरो}}, \emph{pūro}.) Happy? (\textbf{\mw{खुस}}, \emph{khus}.) Timid? (\textbf{\mw{डरपोक}}, \emph{ḍarpok}.) Brave? (\textbf{\mw{बहादुर}}, \emph{bahādur}.) Crazy? (\textbf{\mw{गैलो}}, \emph{gailo}.)
\end{tcolorbox}
